\begin{tikzpicture}[x=1cm,y=1cm,font=\sffamily\footnotesize,
 panel/.style={draw=#1,line width=0.7pt,rounded corners=2pt,fill=#1!6,minimum width=4.6cm,minimum height=10.65cm,inner sep=0pt},
 hdr/.style={font=\sffamily\bfseries\small,text=#1},
 lab/.style={font=\sffamily\scriptsize,text=Slate,align=left,text width=4.15cm},
 val/.style={font=\sffamily\scriptsize,text=Ink,align=left,text width=4.15cm},
 flow/.style={-{Stealth[length=1.6mm]},line width=0.6pt,draw=#1},
 link/.style={{Stealth[length=1.8mm]}-{Stealth[length=1.8mm]},line width=0.7pt,draw=Ink},
 iface/.style={font=\sffamily\tiny,text=Ink,align=center}]
\definecolor{CA}{HTML}{173E60}\definecolor{CW}{HTML}{9A6B00}\definecolor{CI}{HTML}{26714D}
\newcommand{\Harp}{-{Straight Barb[left,length=1.7mm,width=1.4mm]}}
\foreach \n/\x/\c in {a/0/CA,w/5.2/CW,i/10.4/CI}{\node[panel=\c,anchor=north west] (\n) at (\x,0) {};}
\node[hdr=CA,anchor=north west] at (0.22,-0.2) {About DNA};
\node[hdr=CW,anchor=north west] at (5.42,-0.2) {With DNA};
\node[hdr=CI,anchor=north west] at (10.62,-0.2) {Inspired by DNA};
\node[lab,anchor=north west] at (0.22,-0.68) {DNA is \textbf{data}: recorded sequence read by a model};
\node[lab,anchor=north west] at (5.42,-0.68) {DNA is the \textbf{machine}: molecules and reactions compute};
\node[lab,anchor=north west] at (10.62,-0.68) {DNA is a \textbf{design hypothesis}: borrowed organization};
% interfaces across the gaps, labelled above the panels
\draw[link] (4.55,-0.45) -- (5.25,-0.45);
\node[iface,anchor=south] at (4.9,0.06) {sequencing and synthesis};
\draw[link] (9.75,-0.45) -- (10.45,-0.45);
\node[iface,anchor=south] at (10.1,0.06) {typed molecular backend (hypothesis)};
% sketch 1: sequence -> model -> track
\begin{scope}[shift={(0.3,-2.05)}]
\node[draw=CA,fill=white,font=\ttfamily\scriptsize,minimum width=3.9cm,minimum height=0.45cm,anchor=west] at (0,0) {ACGTTGCAACGG\,\dots};
\foreach \k in {0,1,2}{\draw[CA,fill=CA!15] (0.75+\k*1.0,-0.85) rectangle ++(0.8,-0.5);}
\node[font=\sffamily\tiny,text=CA,anchor=north] at (1.95,-1.38) {sequence model};
\draw[flow=CA] (1.95,-0.25) -- (1.95,-0.82);
\draw[flow=CA] (1.95,-1.65) -- (1.95,-2.0);
\foreach \h [count=\k] in {0.15,0.35,0.6,0.3,0.12,0.5,0.75,0.25,0.1}{\fill[CA!70] (0.2+0.38*\k,-2.75) rectangle ++(0.26,\h);}
\draw[Slate,line width=0.4pt] (0.4,-2.75) -- (3.9,-2.75);
\node[font=\sffamily\tiny,text=Slate,anchor=north] at (2.1,-2.8) {predicted track along the sequence};
\end{scope}
% sketch 2: toehold-mediated strand displacement -> readout
\begin{scope}[shift={(5.5,-2.05)}]
\draw[CW,line width=1.4pt,\Harp] (0.2,0) -- (2.6,0);
\draw[CW!60!black,line width=1.4pt,\Harp] (2.6,-0.18) -- (0.6,-0.18);
\foreach \t in {0.7,0.95,...,2.5}{\draw[Slate,line width=0.3pt] (\t,-0.03) -- (\t,-0.15);}
\node[font=\sffamily\tiny,text=Slate,anchor=north] at (0.4,-0.24) {toehold};
\draw[Red,line width=1.4pt,\Harp] (4.1,-0.55) -- (2.9,-0.55);
\node[font=\sffamily\tiny,text=Red,anchor=south] at (3.5,-0.5) {input strand};
\draw[flow=CW] (1.6,-0.45) -- (1.6,-1.15);
\node[font=\sffamily\tiny,text=Slate,anchor=west,align=left] at (1.68,-0.8) {strand displacement};
\draw[CW,line width=1.4pt,\Harp] (0.2,-1.45) -- (2.6,-1.45);
\draw[Red,line width=1.4pt,\Harp] (2.6,-1.63) -- (0.2,-1.63);
\draw[CW!60!black,line width=1.4pt,\Harp] (4.0,-1.55) -- (2.9,-1.55);
\node[star,star points=6,fill=Gold,draw=CW,inner sep=1.4pt] at (4.15,-1.55) {};
\draw[flow=CW] (3.45,-1.75) -- (3.45,-2.2);
\node[font=\sffamily\tiny,text=Slate,anchor=north,align=center] at (2.05,-2.3) {released reporter $\Rightarrow$ fluorescence readout\\[-1pt]half-arrowheads mark $3'$ ends};
\end{scope}
% sketch 3: catalog -> regulated plan -> test
\begin{scope}[shift={(10.6,-2.05)}]
\node[font=\sffamily\tiny,text=Slate,anchor=west] at (-0.05,0.0) {stored catalog};
\foreach \k/\t in {0/clip,1/center,2/prefix,3/scale}{\node[draw=CI,fill=white,font=\sffamily\tiny,minimum width=0.85cm,minimum height=0.36cm] (m\k) at (0.45+\k*0.98,-0.45) {\t};}
\node[draw=CI,fill=CI!15,font=\sffamily\tiny,minimum width=1.6cm,minimum height=0.4cm] (rg) at (1.95,-1.3) {regulator$(c)$};
\draw[flow=CI] (m0.south) -- (rg); \draw[flow=CI] (m2.south) -- (rg);
\node[draw=CI,fill=white,font=\sffamily\tiny,minimum width=2.9cm,minimum height=0.4cm] (pl) at (1.95,-2.05) {expressed plan: clip $\to$ prefix};
\draw[flow=CI] (rg) -- (pl);
\node[draw=CI,dashed,fill=white,font=\sffamily\tiny,minimum width=3.4cm,minimum height=0.4cm] (ab) at (1.95,-2.8) {ablation vs matched baseline};
\draw[flow=CI] (pl) -- (ab);
\end{scope}
% rows: success, evidence, examples
\foreach \x/\a/\b/\c in {
 0.22/{Held-out prediction accuracy and calibration}/{Homology-aware splits; distal-context and mutation interventions}/{DNABERT, Enformer, Nucleotide Transformer, HyenaDNA, Caduceus, Evo, AlphaGenome; Hermon DNA},
 5.42/{The intended output molecules form, at measured yield and speed}/{Measured kinetics and yields with leak, sham and replicate controls}/{Strand-displacement circuits, DNA neural networks, self-assembly, DNA storage},
 10.62/{A measurable advantage that disappears when the borrowed organization is removed}/{Matched-budget ablations and transplants, with a falsifier stated first}/{DOGMA; Evolutor runtime and governance}}{
 \node[lab,anchor=north west] at (\x,-5.4) {\textbf{A successful computation}};
 \node[val,anchor=north west] at (\x,-5.78) {\a};
 \node[lab,anchor=north west] at (\x,-7.0) {\textbf{Strongest evidence}};
 \node[val,anchor=north west] at (\x,-7.38) {\b};
 \node[lab,anchor=north west] at (\x,-8.6) {\textbf{Examples}};
 \node[val,anchor=north west] at (\x,-8.98) {\c};}
% about <-> inspired
\draw[link] (2.3,-10.65) -- ++(0,-0.4) -| (12.7,-10.65);
\node[font=\sffamily\scriptsize,text=Ink,fill=white,inner sep=2pt] at (7.5,-11.05) {model directions and design hypotheses: DOGMA and Hermon DNA are tested as models of DNA \emph{data}};
\end{tikzpicture}
